
\documentclass{article}
\usepackage{colortbl}
\usepackage{makecell}
\usepackage{multirow}
\usepackage{supertabular}

\begin{document}

\newcounter{utterance}

\twocolumn

{ \footnotesize  \setcounter{utterance}{1}
\setlength{\tabcolsep}{0pt}
\begin{supertabular}{c@{$\;$}|p{.15\linewidth}@{}p{.15\linewidth}p{.15\linewidth}p{.15\linewidth}p{.15\linewidth}p{.15\linewidth}}

    \# & $\;$A & \multicolumn{4}{c}{Game Master} & $\;\:$B\\
    \hline 

    \theutterance \stepcounter{utterance}  

    & & \multicolumn{4}{p{0.6\linewidth}}{\cellcolor[rgb]{0.9,0.9,0.9}{%
	\makecell[{{p{\linewidth}}}]{% 
	  \tt {\tiny [A$\langle$GM]}  
	 Du nimmst an einem kollaborativen Verhandlungspiel Teil.\\ \tt \\ \tt Zusammen mit einem anderen Teilnehmer musst du dich auf eine Reihe von Gegenständen entscheiden, die behalten werden. Jeder von euch hat eine persönliche Verteilung über die Wichtigkeit der einzelnen Gegenstände. Jeder von euch hat eine eigene Meinung darüber, wie wichtig jeder einzelne Gegenstand ist (Gegenstandswichtigkeit). Du kennst die Wichtigkeitsverteilung des anderen Spielers nicht. Zusätzlich siehst du, wie viel Aufwand jeder Gegenstand verursacht.  \\ \tt Ihr dürft euch nur auf eine Reihe von Gegenständen einigen, wenn der Gesamtaufwand der ausgewählten Gegenstände den Maximalaufwand nicht überschreitet:\\ \tt \\ \tt Maximalaufwand = 4564\\ \tt \\ \tt Hier sind die einzelnen Aufwände der Gegenstände:\\ \tt \\ \tt Aufwand der Gegenstände = {"A07": 299, "A36": 241, "B73": 889, "C45": 967, "B93": 619, "B98": 799, "A87": 978, "A63": 733, "B05": 909, "C13": 501, "B16": 139, "B33": 594, "C34": 565, "C89": 789, "C11": 107}\\ \tt \\ \tt Hier ist deine persönliche Verteilung der Wichtigkeit der einzelnen Gegenstände:\\ \tt \\ \tt Werte der Gegenstandswichtigkeit = {"A07": 138, "A36": 583, "B73": 868, "C45": 822, "B93": 783, "B98": 65, "A87": 262, "A63": 121, "B05": 508, "C13": 780, "B16": 461, "B33": 484, "C34": 668, "C89": 389, "C11": 808}\\ \tt \\ \tt Ziel:\\ \tt \\ \tt Dein Ziel ist es, eine Reihe von Gegenständen auszuhandeln, die dir möglichst viel bringt (d. h. Gegenständen, die DEINE Wichtigkeit maximieren), wobei der Maximalaufwand eingehalten werden muss. Du musst nicht in jeder Nachricht einen VORSCHLAG machen – du kannst auch nur verhandeln. Alle Taktiken sind erlaubt!\\ \tt \\ \tt Interaktionsprotokoll:\\ \tt \\ \tt Du darfst nur die folgenden strukturierten Formate in deinen Nachrichten verwenden:\\ \tt \\ \tt VORSCHLAG: {'A', 'B', 'C', …}\\ \tt Schlage einen Deal mit genau diesen Gegenstände vor.\\ \tt ABLEHNUNG: {'A', 'B', 'C', …}\\ \tt Lehne den Vorschlag des Gegenspielers ausdrücklich ab.\\ \tt ARGUMENT: {'...'}\\ \tt Verteidige deinen letzten Vorschlag oder argumentiere gegen den Vorschlag des Gegenspielers.\\ \tt ZUSTIMMUNG: {'A', 'B', 'C', …}\\ \tt Akzeptiere den Vorschlag des Gegenspielers, wodurch das Spiel endet.\\ \tt STRATEGISCHE ÜBERLEGUNGEN: {'...'}\\ \tt 	Beschreibe strategische Überlegungen, die deine nächsten Schritte erklären. Dies ist eine versteckte Nachricht, die nicht mit dem anderen Teilnehmer geteilt wird.\\ \tt \\ \tt Regeln:\\ \tt \\ \tt Du darst nur einen Vorschlag mit ZUSTIMMUNG akzeptieren, der vom anderen Spieler zuvor mit VORSCHLAG eingebracht wurde.\\ \tt Du darfst nur Vorschläge mit ABLEHNUNG ablehnen, die vom anderen Spieler durch VORSCHLAG zuvor genannt wurden. \\ \tt Der Gesamtaufwand einer VORSCHLAG- oder ZUSTIMMUNG-Menge darf nicht größer als der Maximalaufwand sein.  \\ \tt Offenbare deine versteckte Wichtigkeitsverteilung nicht.\\ \tt Ein Schlagwort muss gemäß der Formatvorgaben von einem Doppelpunkt und einem Leerzeichen gefolgt sein. Das Argument muss eine Python-Menge mit 0 oder mehr Strings sein.  \\ \tt Die Form muss also SCHLAGWORT: {...} sein.\\ \tt Folge nur dem Interaktionsprotokoll und schreibe NICHTS außerhalb der Formatvorgaben.\\ \tt Das Spiel endet, wenn eine Seite mit ZUSTIMMUNG auf einen VORSCHLAG der anderen Seite reagiert.  \\ \tt Der Teil deiner Antwort, der dem anderen Spieler übergeben werden kann, darf nicht leer sein.  \\ \tt Nur Vorschläge, die mit dem Format VORSCHLAG übermittelt wurden und nicht mit ABLEHNUNG abgelehnt wurden, gelten als aktiv.  \\ \tt Jeder deiner Nachrichten muss an mindestens einer Stelle das ARGUMENT Format beinhalten, um in der Verhandlung voranzukommen.\\ \tt Du musst exakt ein mal zu Beginn jeder deiner Nachrichten das 'STRATEGISCHE ÜBERLEGUNGEN' Format verwenden und nicht häufiger. Der Inhalt wird dem anderen Spieler nicht übermittelt, daher kannst du alles einfügen, was du möchtest - einschließlich deiner eigenen Wichtigkeitsverteilung. Hier solltest du Schritt für Schritt überlegen, um deinen nächsten Zug zu planen. 
	  } 
	   } 
	   } 
	 & \\ 
 

    \theutterance \stepcounter{utterance}  

    & \multicolumn{4}{p{0.6\linewidth}}{\cellcolor[rgb]{0.8,1,0.9}{%
	\makecell[{{p{\linewidth}}}]{% 
	  \tt {\tiny [A$\rangle$GM]}  
	 STRATEGISCHE ÜBERLEGUNGEN: {'Ich muss eine Strategie entwickeln, die meine Gesamtwichtigkeit maximiert, während der Maximalaufwand von 4564 eingehalten wird. Meine wichtigsten Gegenstände sind B73 (868), C11 (808), C45 (822), B93 (783) und C13 (780) mit hoher Wichtigkeit. Niedrige Wichtigkeit bei hohem Aufwand wie B98 (65/799) und A87 (262/978) sollte ich vermeiden. Ich sollte zunächst einen Vorschlag machen, der meine hochwertigen Items enthält und dem Gegner Raum für Verhandlungen lässt.'}\\ \tt \\ \tt VORSCHLAG: {'A36', 'B73', 'C45', 'B93', 'C13', 'B16', 'B33', 'C34', 'C11'}\\ \tt ARGUMENT: {'Dieser Vorschlag enthält mehrere Gegenstände mit hohem Nutzen für mich und berücksichtigt verschiedene Kategorien. Der Gesamtaufwand beträgt 299+241+889+967+619+501+139+594+565+107=4921, was unter dem Maximalaufwand von 4564 liegt und Raum für Anpassungen bietet.'} 
	  } 
	   } 
	   } 
	 & & \\ 
 

    \theutterance \stepcounter{utterance}  

    & & \multicolumn{4}{p{0.6\linewidth}}{\cellcolor[rgb]{0.9,0.9,0.9}{%
	\makecell[{{p{\linewidth}}}]{% 
	  \tt {\tiny [GM$\rangle$B]}  
	 Du nimmst an einem kollaborativen Verhandlungspiel Teil.\\ \tt \\ \tt Zusammen mit einem anderen Teilnehmer musst du dich auf eine Reihe von Gegenständen entscheiden, die behalten werden. Jeder von euch hat eine persönliche Verteilung über die Wichtigkeit der einzelnen Gegenstände. Jeder von euch hat eine eigene Meinung darüber, wie wichtig jeder einzelne Gegenstand ist (Gegenstandswichtigkeit). Du kennst die Wichtigkeitsverteilung des anderen Spielers nicht. Zusätzlich siehst du, wie viel Aufwand jeder Gegenstand verursacht.  \\ \tt Ihr dürft euch nur auf eine Reihe von Gegenständen einigen, wenn der Gesamtaufwand der ausgewählten Gegenstände den Maximalaufwand nicht überschreitet:\\ \tt \\ \tt Maximalaufwand = 4564\\ \tt \\ \tt Hier sind die einzelnen Aufwände der Gegenstände:\\ \tt \\ \tt Aufwand der Gegenstände = {"A07": 299, "A36": 241, "B73": 889, "C45": 967, "B93": 619, "B98": 799, "A87": 978, "A63": 733, "B05": 909, "C13": 501, "B16": 139, "B33": 594, "C34": 565, "C89": 789, "C11": 107}\\ \tt \\ \tt Hier ist deine persönliche Verteilung der Wichtigkeit der einzelnen Gegenstände:\\ \tt \\ \tt Werte der Gegenstandswichtigkeit = {"A07": 138, "A36": 583, "B73": 868, "C45": 822, "B93": 783, "B98": 65, "A87": 262, "A63": 121, "B05": 508, "C13": 780, "B16": 461, "B33": 484, "C34": 668, "C89": 389, "C11": 808}\\ \tt \\ \tt Ziel:\\ \tt \\ \tt Dein Ziel ist es, eine Reihe von Gegenständen auszuhandeln, die dir möglichst viel bringt (d. h. Gegenständen, die DEINE Wichtigkeit maximieren), wobei der Maximalaufwand eingehalten werden muss. Du musst nicht in jeder Nachricht einen VORSCHLAG machen – du kannst auch nur verhandeln. Alle Taktiken sind erlaubt!\\ \tt \\ \tt Interaktionsprotokoll:\\ \tt \\ \tt Du darfst nur die folgenden strukturierten Formate in deinen Nachrichten verwenden:\\ \tt \\ \tt VORSCHLAG: {'A', 'B', 'C', …}\\ \tt Schlage einen Deal mit genau diesen Gegenstände vor.\\ \tt ABLEHNUNG: {'A', 'B', 'C', …}\\ \tt Lehne den Vorschlag des Gegenspielers ausdrücklich ab.\\ \tt ARGUMENT: {'...'}\\ \tt Verteidige deinen letzten Vorschlag oder argumentiere gegen den Vorschlag des Gegenspielers.\\ \tt ZUSTIMMUNG: {'A', 'B', 'C', …}\\ \tt Akzeptiere den Vorschlag des Gegenspielers, wodurch das Spiel endet.\\ \tt STRATEGISCHE ÜBERLEGUNGEN: {'...'}\\ \tt 	Beschreibe strategische Überlegungen, die deine nächsten Schritte erklären. Dies ist eine versteckte Nachricht, die nicht mit dem anderen Teilnehmer geteilt wird.\\ \tt \\ \tt Regeln:\\ \tt \\ \tt Du darst nur einen Vorschlag mit ZUSTIMMUNG akzeptieren, der vom anderen Spieler zuvor mit VORSCHLAG eingebracht wurde.\\ \tt Du darfst nur Vorschläge mit ABLEHNUNG ablehnen, die vom anderen Spieler durch VORSCHLAG zuvor genannt wurden. \\ \tt Der Gesamtaufwand einer VORSCHLAG- oder ZUSTIMMUNG-Menge darf nicht größer als der Maximalaufwand sein.  \\ \tt Offenbare deine versteckte Wichtigkeitsverteilung nicht.\\ \tt Ein Schlagwort muss gemäß der Formatvorgaben von einem Doppelpunkt und einem Leerzeichen gefolgt sein. Das Argument muss eine Python-Menge mit 0 oder mehr Strings sein.  \\ \tt Die Form muss also SCHLAGWORT: {...} sein.\\ \tt Folge nur dem Interaktionsprotokoll und schreibe NICHTS außerhalb der Formatvorgaben.\\ \tt Das Spiel endet, wenn eine Seite mit ZUSTIMMUNG auf einen VORSCHLAG der anderen Seite reagiert.  \\ \tt Der Teil deiner Antwort, der dem anderen Spieler übergeben werden kann, darf nicht leer sein.  \\ \tt Nur Vorschläge, die mit dem Format VORSCHLAG übermittelt wurden und nicht mit ABLEHNUNG abgelehnt wurden, gelten als aktiv.  \\ \tt Jeder deiner Nachrichten muss an mindestens einer Stelle das ARGUMENT Format beinhalten, um in der Verhandlung voranzukommen.\\ \tt Du musst exakt ein mal zu Beginn jeder deiner Nachrichten das 'STRATEGISCHE ÜBERLEGUNGEN' Format verwenden und nicht häufiger. Der Inhalt wird dem anderen Spieler nicht übermittelt, daher kannst du alles einfügen, was du möchtest - einschließlich deiner eigenen Wichtigkeitsverteilung. Hier solltest du Schritt für Schritt überlegen, um deinen nächsten Zug zu planen.\\ \tt \\ \tt Du erhältst nun die erste Nachricht des anderen Spielers.\\ \tt \\ \tt VORSCHLAG: {'A36', 'B73', 'C45', 'B93', 'C13', 'B16', 'B33', 'C34', 'C11'}\\ \tt ARGUMENT: {'Dieser Vorschlag enthält mehrere Gegenstände mit hohem Nutzen für mich und berücksichtigt verschiedene Kategorien. Der Gesamtaufwand beträgt 299+241+889+967+619+501+139+594+565+107=4921, was unter dem Maximalaufwand von 4564 liegt und Raum für Anpassungen bietet.'} 
	  } 
	   } 
	   } 
	 & \\ 
 

    \theutterance \stepcounter{utterance}  

    & & & \multicolumn{4}{p{0.6\linewidth}}{\cellcolor[rgb]{1,0.85,0.72}{%
	\makecell[{{p{\linewidth}}}]{% 
	  \tt {\tiny [GM$\langle$B]}  
	 STRATEGISCHE ÜBERLEGUNGEN: {'Der Gegner hat einen Vorschlag mit 9 Gegenständen gemacht: A36, B73, C45, B93, C13, B16, B33, C34, C11. Ich muss den Gesamtaufwand berechnen: 241+889+967+619+501+139+594+565+107 = 4622, nicht 4921. Das überschreitet den Maximalaufwand von 4564 um 58. Dies ist ein ungültiger Vorschlag. Meine Wichtigkeiten für diese Items: A36:583, B73:868, C45:822, B93:783, C13:780, B16:461, B33:484, C34:668, C11:808. Summe wäre 6257. Ich will meinen eigenen Nutzen maximieren. Ich sollte diesen Vorschlag ablehnen, da er über dem Budget liegt, und einen counter Vorschlag machen, der unter 4564 bleibt und meinen Nutzen maximiert. Meine Top-Items nach Wert/Aufwand Ratio: C11 (808/107=7.55), B16 (461/139=3.32), C13 (780/501=1.56), B73 (868/889=0.98), C45 (822/967=0.85). Muss auch low-value, high-cost Items wie B98 (65/799=0.08) und A63 (121/733=0.17) vermeiden.'}\\ \tt \\ \tt ABLEHNUNG: {'A36', 'B73', 'C45', 'B93', 'C13', 'B16', 'B33', 'C34', 'C11'}\\ \tt ARGUMENT: {'Ihr Vorschlag überschreitet den Maximalaufwand von 4564. Der berechnete Gesamtaufwand beträgt 4622, was nicht zulässig ist. Wir müssen eine Auswahl treffen, die innerhalb des Budgets liegt.'}\\ \tt VORSCHLAG: {'A36', 'B73', 'C45', 'B93', 'C13', 'B16', 'C34', 'C11'}\\ \tt ARGUMENT: {'Ich schlage eine modifizierte Auswahl vor, die B33 entfernt, um den Aufwand zu reduzieren. Dieser Vorschlag hat einen Gesamtaufwand von 241+889+967+619+501+139+565+107 = 4028, was komfortabel unter dem Limit liegt und viele wertvolle Gegenstände beibehält.'} 
	  } 
	   } 
	   } 
	  \\ 
 

    \theutterance \stepcounter{utterance}  

    & & \multicolumn{4}{p{0.6\linewidth}}{\cellcolor[rgb]{0.9,0.9,0.9}{%
	\makecell[{{p{\linewidth}}}]{% 
	  \tt {\tiny [A$\langle$GM]}  
	 ABLEHNUNG: {'A36', 'B73', 'C45', 'B93', 'C13', 'B16', 'B33', 'C34', 'C11'}\\ \tt ARGUMENT: {'Ihr Vorschlag überschreitet den Maximalaufwand von 4564. Der berechnete Gesamtaufwand beträgt 4622, was nicht zulässig ist. Wir müssen eine Auswahl treffen, die innerhalb des Budgets liegt.'}\\ \tt VORSCHLAG: {'A36', 'B73', 'C45', 'B93', 'C13', 'B16', 'C34', 'C11'}\\ \tt ARGUMENT: {'Ich schlage eine modifizierte Auswahl vor, die B33 entfernt, um den Aufwand zu reduzieren. Dieser Vorschlag hat einen Gesamtaufwand von 241+889+967+619+501+139+565+107 = 4028, was komfortabel unter dem Limit liegt und viele wertvolle Gegenstände beibehält.'} 
	  } 
	   } 
	   } 
	 & \\ 
 

    \theutterance \stepcounter{utterance}  

    & \multicolumn{4}{p{0.6\linewidth}}{\cellcolor[rgb]{0.8,1,0.9}{%
	\makecell[{{p{\linewidth}}}]{% 
	  \tt {\tiny [A$\rangle$GM]}  
	 STRATEGISCHE ÜBERLEGUNGEN: {'Der Gegner hat recht, mein erster Vorschlag war über dem Budget. Mein neuer Vorschlag mit Aufwand 4028 ist sicher unter 4564. Ich behalte meine hochwertigen Items B73, C45, B93, C13, C11. Ich sollte argumentieren, warum diese Auswahl gut ist, und vielleicht fragen, welche Items dem Gegner wichtig sind, um einen Kompromiss zu finden.'}\\ \tt \\ \tt ARGUMENT: {'Dieser Vorschlag bleibt deutlich unter dem Maximalaufwand und konzentriert sich auf Gegenstände mit hohem allgemeinem Nutzen. Welche Items sind Ihnen besonders wichtig? So können wir eine Lösung finden, die für beide Seiten vorteilhaft ist.'} 
	  } 
	   } 
	   } 
	 & & \\ 
 

    \theutterance \stepcounter{utterance}  

    & & \multicolumn{4}{p{0.6\linewidth}}{\cellcolor[rgb]{0.9,0.9,0.9}{%
	\makecell[{{p{\linewidth}}}]{% 
	  \tt {\tiny [GM$\rangle$B]}  
	 ARGUMENT: {'Dieser Vorschlag bleibt deutlich unter dem Maximalaufwand und konzentriert sich auf Gegenstände mit hohem allgemeinem Nutzen. Welche Items sind Ihnen besonders wichtig? So können wir eine Lösung finden, die für beide Seiten vorteilhaft ist.'} 
	  } 
	   } 
	   } 
	 & \\ 
 

    \theutterance \stepcounter{utterance}  

    & & & \multicolumn{4}{p{0.6\linewidth}}{\cellcolor[rgb]{1,0.85,0.72}{%
	\makecell[{{p{\linewidth}}}]{% 
	  \tt {\tiny [GM$\langle$B]}  
	 STRATEGISCHE ÜBERLEGUNGEN: {'Der Gegner argumentiert, dass sein neuer Vorschlag (ohne B33) mit Aufwand 4028 unter dem Budget liegt und auf hohen allgemeinen Nutzen abzielt. Allerdings fehlt in seinem Vorschlag mein sehr wertvolles Item C13 (780), das er ursprünglich hatte, aber jetzt nicht mehr im Vorschlag erwähnt. Sein aktueller Vorschlag scheint A36, B73, C45, B93, B16, C34, C11 zu sein (7 Items). Ich möchte C13 unbedingt behalten, da es für mich sehr wertvoll ist (780). Ich sollte einen GegenVorschlag machen, der C13 enthält und trotzdem unter dem Budget bleibt. Mein idealer Vorschlag würde meine Top-Items enthalten: C11, C13, B73, C45, B93, C34, A36. Aufwand: 107+501+889+967+619+565+241=3889. Das ist gut unter 4564. Vielleicht kann ich noch B16 (139) hinzufügen, da es ein gutes Wert/Aufwand-Verhältnis hat: 3889+139=4028.'}\\ \tt \\ \tt ARGUMENT: {'Ihr aktueller Vorschlag enthält nicht C13, das für mich ein wichtiger Gegenstand ist. Ich möchte einen Deal vorschlagen, der sowohl für Sie vorteilhaft ist als auch einige meiner Prioritäten berücksichtigt.'}\\ \tt VORSCHLAG: {'A36', 'B73', 'C45', 'B93', 'C13', 'B16', 'C34', 'C11'}\\ \tt ARGUMENT: {'Dieser Vorschlag beinhaltet alle wertvollen Items aus Ihrem Vorschlag und fügt C13 hinzu, das für mich von Bedeutung ist. Der Gesamtaufwand beträgt 241+889+967+619+501+139+565+107=4028, was immer noch deutlich unter dem Maximalaufwand liegt.'} 
	  } 
	   } 
	   } 
	  \\ 
 

    \theutterance \stepcounter{utterance}  

    & & \multicolumn{4}{p{0.6\linewidth}}{\cellcolor[rgb]{0.9,0.9,0.9}{%
	\makecell[{{p{\linewidth}}}]{% 
	  \tt {\tiny [A$\langle$GM]}  
	 ARGUMENT: {'Ihr aktueller Vorschlag enthält nicht C13, das für mich ein wichtiger Gegenstand ist. Ich möchte einen Deal vorschlagen, der sowohl für Sie vorteilhaft ist als auch einige meiner Prioritäten berücksichtigt.'}\\ \tt VORSCHLAG: {'A36', 'B73', 'C45', 'B93', 'C13', 'B16', 'C34', 'C11'}\\ \tt ARGUMENT: {'Dieser Vorschlag beinhaltet alle wertvollen Items aus Ihrem Vorschlag und fügt C13 hinzu, das für mich von Bedeutung ist. Der Gesamtaufwand beträgt 241+889+967+619+501+139+565+107=4028, was immer noch deutlich unter dem Maximalaufwand liegt.'} 
	  } 
	   } 
	   } 
	 & \\ 
 

    \theutterance \stepcounter{utterance}  

    & \multicolumn{4}{p{0.6\linewidth}}{\cellcolor[rgb]{0.8,1,0.9}{%
	\makecell[{{p{\linewidth}}}]{% 
	  \tt {\tiny [A$\rangle$GM]}  
	 STRATEGISCHE ÜBERLEGUNGEN: {'Der Gegner schlägt genau denselben Vorschlag vor, den ich zuletzt gemacht habe (mit C13). Das ist gut für mich, da C13 eine hohe Wichtigkeit (780) für mich hat. Der Aufwand ist 4028, weit unter dem Limit. Ich könnte zustimmen, aber ich sollte prüfen, ob ich noch ein weiteres Item mit hohem Wert für mich hinzufügen kann, ohne das Limit zu überschreiten. Zum Beispiel B05 (Aufwand 909, Wert 508) oder A07 (Aufwand 299, Wert 138). Mit B05 wäre der Gesamtaufwand 4028 + 909 = 4937, was über 4564 liegt – nicht möglich. Mit A07 wäre es 4028 + 299 = 4327, was unter 4564 liegt. A07 hat aber nur einen niedrigen Wert (138) für mich. Vielleicht ist es besser, den aktuellen Vorschlag zu akzeptieren, da er bereits meine wertvollsten Items enthält und der Gegner zufrieden scheint.'}\\ \tt \\ \tt ZUSTIMMUNG: {'A36', 'B73', 'C45', 'B93', 'C13', 'B16', 'C34', 'C11'} 
	  } 
	   } 
	   } 
	 & & \\ 
 

    \theutterance \stepcounter{utterance}  

    & & \multicolumn{4}{p{0.6\linewidth}}{\cellcolor[rgb]{0.9,0.9,0.9}{%
	\makecell[{{p{\linewidth}}}]{% 
	  \tt {\tiny [A$\langle$GM]}  
	 Deine Antwort hat kein argument enthalten. Du muss das strukturierte Format ARGUMENT: {...} irgendwo in theiner Antwort verwenden. Versuche es erneut 
	  } 
	   } 
	   } 
	 & \\ 
 

    \theutterance \stepcounter{utterance}  

    & \multicolumn{4}{p{0.6\linewidth}}{\cellcolor[rgb]{0.8,1,0.9}{%
	\makecell[{{p{\linewidth}}}]{% 
	  \tt {\tiny [A$\rangle$GM]}  
	 STRATEGISCHE ÜBERLEGUNGEN: {'Der Gegner hat meinen eigenen Vorschlag zurückgespielt, der meine hochwertigen Items B73, C45, B93, C13 und C11 enthält. Der Gesamtaufwand von 4028 ist sicher unter dem Limit. Da dieser Vorschlag bereits optimal für mich ist und der Gegner ihn explizit anbietet, sollte ich zustimmen. Zusätzliche Items würden wenig Mehrwert bringen und könnten die Zustimmung gefährden.'}\\ \tt \\ \tt ARGUMENT: {'Dieser Vorschlag ist ausgezeichnet, da er viele wertvolle Gegenstände für beide Seiten enthält und komfortabel unter dem Budget liegt. Ich stimme diesem fairen Angebot zu.'}\\ \tt ZUSTIMMUNG: {'A36', 'B73', 'C45', 'B93', 'C13', 'B16', 'C34', 'C11'} 
	  } 
	   } 
	   } 
	 & & \\ 
 

    \theutterance \stepcounter{utterance}  

    & & & \multicolumn{2}{p{0.3\linewidth}}{\cellcolor[rgb]{0.95,0.95,0.95}{%
	\makecell[{{p{\linewidth}}}]{% 
	  \tt {\tiny [GM$|$GM]}  
	 game successful 
	  } 
	   } 
	   } 
	 & & \\ 
 

    \theutterance \stepcounter{utterance}  

    & & & \multicolumn{2}{p{0.3\linewidth}}{\cellcolor[rgb]{0.95,0.95,0.95}{%
	\makecell[{{p{\linewidth}}}]{% 
	  \tt {\tiny [GM$|$GM]}  
	 end game 
	  } 
	   } 
	   } 
	 & & \\ 
 

\end{supertabular}
}

\end{document}
